\subsection{Central scalar refinements and full peripheral lifts}

\begin{statement}{D.50}{Proved}{Every multiplier-one pronilpotent peripheral source lifts to the full source}
Write
\[
 \mathcal F=\widehat F(x_1,x_2),\qquad x_3=(x_1x_2)^{-1},
 \qquad \mathcal F_{\rm nil}=\mathcal F/R_{\rm nil}(\mathcal F),
 \tag{D.50.1}
\]
where \(R_{\rm nil}(G)\) is the intersection of the kernels of all
continuous homomorphisms from \(G\) to finite nilpotent groups.
For every continuous automorphism \(\alpha\) of
\(\mathcal F_{\rm nil}\) and every specified \(b_i\in\mathcal F_{\rm nil}\)
such that
\[
          \alpha(\bar x_i)=b_i\bar x_i b_i^{-1}\quad(1\leq i\leq3),
 \tag{D.50.2}
\]
there are a continuous automorphism \(\Phi\) of \(\mathcal F\)
and actual lifts \(h_i\in\mathcal F\) of all three \(b_i\) with
\[
 q_{\rm nil}\Phi=\alpha q_{\rm nil},\qquad
                         \Phi(x_i)=h_ix_ih_i^{-1}.
 \tag{D.50.3}
\]
If \(b_1=1\), the lift can be chosen to fix \(x_1\) literally.

There is a corresponding statement for the actual modular source
\[
 \mathcal B=\widehat{C_2*C_3}=\langle s,r\mid s^2=r^3=1\rangle,
 \quad \tau=sr,\quad
 F=\ker(\mathcal B\to\mathfrak S_3),\quad
 \mathcal B_{\rm nil}=\mathcal B/R_{\rm nil}(F).
 \tag{D.50.4}
\]
Every continuous automorphism of \(\mathcal B_{\rm nil}\) preserving
the positive classes of \(s,r,\tau\), with any specified three
conjugators there, lifts to a continuous automorphism of \(\mathcal B\)
with those conjugators lifted literally. If the old map fixes \(\tau\)
and its specified \(\tau\)-conjugator is \(1\), the full lift can be
chosen to fix \(\tau\).

In particular, for every \(\ell\geq5\), every actual multiplier-one
pro-\(\ell\) gentle pair \(f_\ell\in F^{(\ell)\prime}\) whose pure
map is an automorphism has a full profinite gentle lift
\(f\in F'\) with the literal image \(f_\ell\).
Its normalized pure map is
\(x\mapsto x,\ y\mapsto f^{-1}yf\), with
\(x=\tau^{-2}\), \(y=sxs\), and its outer class centralizes the
actual \(\mathfrak S_3\)-action. The full pro-\(\ell\) pair is
retained, not merely a finite graded approximation.

The analogous full-source lifting statements hold for the whole
prosolvable quotients as well. All these assertions follow from
the finite central refinement and complete fibre theorem proved below.
\end{statement}

\begin{proof}
We first count every finite proper fibre, retaining its actual source.
The centrality of the resulting refinement is used explicitly in
the joint-generation and inverse-limit steps.

\paragraph{Finite scalar covers for every subgroup image.}
Let
\[
 1\longrightarrow A\longrightarrow E\xrightarrow{\pi}D\longrightarrow1
 \tag{D.50.5}
\]
be any finite extension with \([E,A]=A\).
The source \(\Lambda\) is either \(\mathcal F\), with the three
peripherals in \textup{(D.50.1)}, or \(\mathcal B\), with
\(x_1=s,x_2=r,x_3=\tau^{-1}\).
Fix an actual epimorphism \(q_E:\Lambda\twoheadrightarrow E\), and put
\(q_0=\pi q_E\), \(X_i=q_E(x_i)\), \(d_i=q_0(x_i)\).
Let \(T_i\) be the \(A\)-conjugacy class of \(X_i\).
More generally the count argument below works with any specified
\(A\)-conjugacy invariant subsets of these three cosets.

Retain every subgroup \(J\leq E\) with \(\pi(J)=D\) and
\(T_i\cap J\ne\varnothing\) for all \(i\).
Put \(M=J\cap A\), so \(J/M=D\), and \(e=\exp(J)\).
For each irreducible complex representation
\(\rho:M\to\operatorname{GL}_d(\mathbf C)\) whose equivalence class
is fixed by conjugation by \(J\), define the actual finite group
\[
 \begin{aligned}
 \widetilde J_\rho=\{(g,B):\;&g\in J,\ B\in\operatorname{GL}_d(\mathbf C),\\
       &B\rho(m)B^{-1}=\rho(gmg^{-1})\ (m\in M),\\
                         &\det B\in\mu_e\},\\
 (g,B)(h,C)&=(gh,BC).
 \end{aligned}
 \tag{D.50.6}
\]
Invariance of \(\rho\) supplies an implementing matrix over every
\(g\); multiplying it by a scalar makes its determinant belong to
\(\mu_e\). Schur's lemma makes all implementing matrices scalar
multiples of one another. Thus the projection onto \(J\) is onto,
with exactly \(de\) elements over each \(g\), and has central kernel
\(\{(1,zI):z\in\mu_{de}\}\). In particular this is a finite group,
with its displayed multiplication.

The map
\(m\mapsto(m,\rho(m))\) embeds \(M\) as a normal subgroup:
\((\det\rho(m))^e=\det\rho(m^e)=1\), and conjugation in
\textup{(D.50.6)} gives exactly the graph of \(gmg^{-1}\).
Its intersection with the scalar kernel is trivial, because its
first coordinate is \(m\). This argument retains representations
with a nontrivial kernel and with a nontrivial determinant.
Consequently
\[
 Q_\rho=\widetilde J_\rho/\{(m,\rho(m)):m\in M\},
 \qquad
 1\longrightarrow\mu_{de}\longrightarrow Q_\rho
                \longrightarrow D\longrightarrow1
 \tag{D.50.7}
\]
is an actual finite central extension. There is an actual isomorphism
\[
 \widetilde J_\rho\longrightarrow J\times_D Q_\rho,
                \qquad (g,B)\longmapsto(g,[(g,B)]).
 \tag{D.50.8}
\]
Its kernel is trivial by the graph intersection just proved.
For surjectivity, lift the prescribed \(Q_\rho\)-coordinate to
\((g',B')\); if the prescribed first coordinate is \(g\), multiply
this lift by the graph element for \(g(g')^{-1}\in M\).
This leaves the \(Q_\rho\)-coordinate unchanged and gives first
coordinate \(g\). Denote by \(\Psi_\rho\) the matrix representation
of this pullback obtained from its \(B\)-coordinate.

\paragraph{Every auxiliary map is a map from the actual source.}
For \(\Lambda=\mathcal F\), choose any lifts of \(d_1,d_2\) to
\(Q_\rho\). They define a continuous homomorphism
\(R^0_\rho:\mathcal F\to Q_\rho\) above \(q_0\); its third
peripheral value is their inverse product.

For \(\Lambda=\mathcal B\), choose
\(a_s\in T_1\cap J\) and \(a_r\in T_2\cap J\).
They have orders dividing two and three respectively.
An implementing matrix for \(a_s\) has scalar square by Schur's
lemma, so scale it to obtain \(B_s^2=I\). If \(a_s=1\), take
\(B_s=I\). In the nontrivial case \(2\mid e\), so
\(\det B_s\in\mu_2\subset\mu_e\).
In the same way choose \(B_r^3=I\); either \(a_r=1\) and
\(B_r=I\), or \(3\mid e\) and \(\det B_r\in\mu_3\subset\mu_e\).
The classes of \((a_s,B_s)\) and \((a_r,B_r)\) in
\(Q_\rho\) therefore define an actual homomorphism
\[
                  R^0_\rho:\mathcal B\longrightarrow Q_\rho
 \tag{D.50.9}
\]
above \(q_0\). No relation is imposed on their product beyond
the defining relation for \(\tau=sr\).
Subgroups omitted because some \(T_i\cap J\) is empty will have
zero count on both sides of the comparison, as proved below.

Take \(D^\dagger\) to be the actual joint image of \(q_0\) and
all these \(R^0_\rho\) in their fibre product over \(D\).
Write its source map as \(q^\dagger\). Then
\[
 q^\dagger:\Lambda\twoheadrightarrow D^\dagger,
 \quad p:D^\dagger\twoheadrightarrow D,
 \quad pq^\dagger=q_0,
 \quad Z:=\ker p\text{ is finite central abelian}.
 \tag{D.50.10}
\]
Indeed the kernel embeds in the product of the central kernels
\(\mu_{de}\) in \textup{(D.50.7)}, and hence is central in the
actual joint image as well. This construction has only finitely
many subgroups and irreducible representations. It retains the
actual quotients and source generators throughout.

\paragraph{The exact pants coefficient formula.}
Here is the finite counting identity needed in the proof.
Fix any finite extension \(M\to J\to D\), a homomorphism
\(\chi:\mathcal F\twoheadrightarrow D\), and subsets
\(U_i\subset J_{\chi(x_i)}\) invariant under conjugation by \(M\).
For each \(J\)-invariant irreducible \(\rho\) of \(M\), choose
any actual lift \(R_\rho:\mathcal F\to Q_\rho\) of \(\chi\),
with \(Q_\rho\) and \(\Psi_\rho\) from
\textup{(D.50.6)--(D.50.8)}. The number of homomorphisms
\(\varphi:\mathcal F\to J\) above \(\chi\) with
\(\varphi(x_i)\in U_i\) is exactly
\[
 \frac1{|M|}\sum_{\substack{\rho\in\operatorname{Irr}(M)\\
                            J\text{ fixes }[\rho]}}
 \frac1{\dim\rho}
       \prod_{i=1}^3
         \left(\sum_{a\in U_i}
            \operatorname{tr}\Psi_\rho(a,R_\rho(x_i))\right).
 \tag{D.50.11}
\]
Every element in a trace is a literal element of the pullback
\(J\times_DQ_\rho\).

To prove the formula, set \(B_i=\sum_{a\in U_i}a\in\mathbf C[J]\).
These sums centralize \(\mathbf C[M]\), and
\(B_1B_2B_3\in\mathbf C[M]\), because the three quotient values
have product one. For an irreducible \(\rho\), let \(e_\rho\)
be its central primitive idempotent in \(\mathbf C[M]\).
For \(a\) in the coset of \(d\in D\), conjugation takes
\(e_\rho\) to the idempotent for the transported representation.
Thus \(e_\rho a e_\rho=0\) unless \(d\) fixes \([\rho]\).
Since \(B_i\) centralizes \(e_\rho\),
\(e_\rho B_i=e_\rho B_ie_\rho\), so this block vanishes when
the \(i\)-th quotient value moves \([\rho]\).
The first two quotient values generate \(D\), and hence only
the \(J\)-invariant irreducibles survive in the product.

For such a \(\rho\), lift a tuple \((a_1,a_2,a_3)\) using
the three prescribed \(R_\rho(x_i)\). Their product has second
coordinate one in \(Q_\rho\), so by the exact pullback
\textup{(D.50.8)} it is the graph element for
\(a_1a_2a_3\in M\). Its matrix is therefore
\(\rho(a_1a_2a_3)\).
Each sum of lifted matrices for one \(U_i\) commutes with
\(\rho(M)\): conjugation by the graph of \(m\) replaces
\(a\) by \(mam^{-1}\) and leaves its \(Q_\rho\)-coordinate
unchanged. Schur's lemma makes this sum the scalar matrix
\((\dim\rho)^{-1}\sum_{a\in U_i}
\operatorname{tr}\Psi_\rho(a,R_\rho(x_i))\).
Finally the regular-character identity
\(\mathbf1_{\{1\}}(m)=|M|^{-1}\sum_\rho
(\dim\rho)\operatorname{tr}\rho(m)\)
extracts the identity coefficient of \(B_1B_2B_3\).
The three scalar denominators, the matrix trace, and this
regular-character factor give exactly \textup{(D.50.11)}.
This proves the required genus-zero formula directly.

For comparison, on an oriented surface of genus \(g\) with
\(b\geq1\) boundaries the same proof gives
\[
 |M|^{2g-1}
 \sum_{\rho\text{ fixed by the quotient image}}
   (\dim\rho)^{2-2g-b}
   \prod_{i=1}^b\sum_{a\in U_i}
                    \operatorname{tr}\Psi_\rho(a,R_\rho(c_i)).
 \tag{D.50.12}
\]
For the exact handle factor, write its matrices as
\(\rho(m)X,\rho(n)Y\). Conjugation averaging followed by
matrix-coefficient orthogonality gives
\[
 \sum_{m,n\in M}[\rho(m)X,\rho(n)Y]
                  =(|M|/\dim\rho)^2I.
 \tag{D.50.13}
\]
Only irreducibles fixed by both handle quotient values can
contribute. To check this without assuming an extension of
a noninvariant representation, put
\(P_{a,b}=\sum_{A\in J_a,B\in J_b}[A,B]\).
It centralizes \(\mathbf C[M]\), and, for a fixed lift \(\widetilde a\),
\[
 P_{a,b}=\sum_{B\in J_b}
    \operatorname{Av}_M(\widetilde aB\widetilde a^{-1})B^{-1},
 \qquad \operatorname{Av}_M(T)=\sum_{m\in M}mTm^{-1}.
 \tag{D.50.14}
\]
Sandwiching this identity by \(e_\rho\) forces \(b\) to fix
\([\rho]\), since the averaged term commutes with \(e_\rho\).
The linear inversion anti-involution of the group algebra sends
\(e_\rho\) to \(e_{\rho^\vee}\) and \(P_{a,b}\) to \(P_{b,a}\),
forcing the same condition for \(a\).
This proves every factor and every vanishing in
\textup{(D.50.12)}. The scalar-cover construction is an explicit
finite model for the projective representation classes of
Turaev, Lemma~2.1 and Section~2.3; the positive-genus section
count is Theorem~1.2 there. Our proof above includes the pants
case used here.\cite{ATuraev2009}

\paragraph{Transporting every subgroup count, including all proper images.}
Let \(q_1^\dagger:\Lambda\twoheadrightarrow D^\dagger\) be any
epimorphism, with specified \(r_i\in D^\dagger\) such that
\(q_1^\dagger(x_i)=r_iq^\dagger(x_i)r_i^{-1}\).
Put \(q_1=pq_1^\dagger\), \(h_i=p(r_i)\), choose any
\(H_i\in E\) above \(h_i\), and set
\(T_i^1=H_iT_iH_i^{-1}\), \(T_i^0=T_i\).
For every \(J\leq E\) onto \(D\), choose \(j_i\in J\) above
\(h_i\). Since \(H_ij_i^{-1}\in A\) and the \(T_i\) are
\(A\)-conjugacy invariant, the exact equality is
\[
                 T_i^1\cap J=j_i(T_i^0\cap J)j_i^{-1}.
 \tag{D.50.15}
\]
Indeed \(j_iT_i^0j_i^{-1}\) is again \(A\)-conjugacy invariant,
as \(j_i\) normalizes \(A\). This proves the equality before
intersecting with \(J\), and conjugation by \(j_i\) preserves \(J\).
In particular, every omitted empty-intersection subgroup has
zero count on both sides.

For a retained \((J,\rho)\), project \(q^\dagger,q_1^\dagger,r_i\)
to \(Q_\rho\), obtaining \(R_\rho^0,R_\rho^1,\widetilde r_i\).
They satisfy the literal equality
\(R_\rho^1(x_i)=\widetilde r_iR_\rho^0(x_i)\widetilde r_i^{-1}\).
The pair \((j_i,\widetilde r_i)\) belongs to
\(J\times_DQ_\rho\). By \textup{(D.50.8)}, its actual matrix
conjugates
\(\Psi_\rho(a,R_\rho^0(x_i))\) to
\(\Psi_\rho(j_i a j_i^{-1},R_\rho^1(x_i))\).
Thus every boundary trace sum in \textup{(D.50.11)} is exactly
the same for the old and new data, with no scalar ambiguity.
The formula gives equality of the full homomorphism counts
into this fixed subgroup \(J\).

Write these counts as \(N_j(J)\), and give count zero to
subgroups not mapping onto \(D\), since \(q_j\) is onto.
Every homomorphism has one actual image subgroup. Hence
ordinary inversion on the finite subgroup lattice gives
\[
 \#\{f:\Lambda\twoheadrightarrow E:\pi f=q_j,
                     \ f(x_i)\in T_i^j\}
       =\sum_{J\leq E}\mu(J,E)N_j(J).
 \tag{D.50.16}
\]
Every summand agrees for \(j=0,1\), proving equality of the
complete proper fibres. In the modular case these are counts
from \(\mathcal B\), because membership in \(T_1^j,T_2^j\)
forces the two generator orders to divide two and three;
every counted pants map therefore factors through the actual
free product. The given \(q_E\) makes the reference count
positive. No proper nontrivial kernel image was discarded.

\paragraph{The exact joint-generation condition.}
The count argument used only normality of \(A\). We now use
\([E,A]=A\). Consider the joint image in
\(E\times_DD^\dagger\) of a counted epimorphism and the
prescribed \(q_1^\dagger\); the reference image is the same
case with \(q_E,q^\dagger\).
On the common base kernel it is a subdirect product of
\(A\) and \(Z=\ker(D^\dagger\to D)\).
Let \(C\) be its common quotient from Goursat's lemma.
Conjugation by the full joint image induces the actual
\(E\)-action on its \(A\)-coordinate, whereas it fixes its
\(Z\)-coordinate because \(Z\) is central in \(D^\dagger\).
Therefore \(E\) acts trivially on \(C\), and the quotient
map from \(A\) kills \([E,A]=A\). Thus \(C=1\), proving
\[
 \operatorname{im}(f,q_1^\dagger)
                        =E\times_DD^\dagger.
 \tag{D.50.17}
\]
The literal correspondence \(f\mapsto(f,q_1^\dagger)\)
therefore identifies all the proper fibres in
\textup{(D.50.16)} with the complete fibres above the
specified \(D^\dagger\)-map, retaining the conjugators
\((H_i,r_i)\) and their \(A\)-corrections.

The full ordinary conjugator union is also accounted for.
The actual image of
\(C_{D^\dagger}(q^\dagger(x_i))\) in \(C_D(d_i)\) acts
on the \(A\)-class of \(X_i\), by any lift to \(E\).
Two lifts differ by \(A\), so this is a well-defined
action on classes. If its orbit size is \(t_i\) and
the count in \textup{(D.50.16)} is \(n_0>0\), the complete
ordinary fibre above every prescribed \(q_1^\dagger\)
has exactly \(t_1t_2t_3n_0\) elements.
Conjugation by \((H_i,r_i)\) identifies the old and new
orbits, every member is obtained by replacing \(r_i\)
by an actual centralizer coset, and distinct triples of
\(A\)-classes have disjoint literal fibres. This proves
the claimed count without enlarging any centralizer image.

\paragraph{Lower central residuals and nilpotent refinements.}
Throughout this proof the lower central subgroups of a profinite
group are closed. They give the exact residual comparison
\[
 R_{\rm nil}(G)=\bigcap_{n\geq1}\gamma_nG,\qquad
 f(R_{\rm nil}(G))=\gamma_\infty(f(G))
       \quad\text{for every finite continuous image }f(G).
\]
Indeed a finite nilpotent quotient kills a lower central term.
Conversely, a nonzero element of a quotient by one closed lower
central term is separated by a finite quotient of that bounded-class
nilpotent profinite group. This proves the first equality.
For the second, the nonempty fibres of a fixed finite image value
in the descending compact subgroups have nonempty intersection;
hence image commutes with this intersection.
For a finite group \(G\), its descending lower central
series stabilizes. Denote its stable value by
\(\gamma_\infty(G)\). Directly from the definition,
\[
 [G,\gamma_\infty(G)]=\gamma_\infty(G),\qquad
 G/\gamma_\infty(G)\text{ is nilpotent}.
 \tag{D.50.18}
\]
Every nilpotent quotient kills this subgroup: a quotient
of class \(c\) kills \(\gamma_{c+1}(G)\), which contains
the stable value. Thus the displayed quotient is the
actual maximal nilpotent quotient of \(G\).
A central extension of a group of class \(c\) has class
at most \(c+1\), because its \((c+1)\)-st lower central
subgroup lies in the central kernel. Consequently, taking
\(A=\gamma_\infty(E)\) in \textup{(D.50.5)} makes both
\(D\) and the actual \(D^\dagger\) finite nilpotent groups.
Every auxiliary map therefore factors through
\(\mathcal F_{\rm nil}\).

For the modular source, include its specified
\(\mathfrak S_3\) quotient in every finite quotient
\(E\). Then \(N=q_E(F)=\ker(E\to\mathfrak S_3)\).
Take \(A=\gamma_\infty(N)\). It is characteristic in
\(N\), hence normal in \(E\), and
\([E,A]=A\), since already \([N,A]=A\).
The actual pure subgroup of \(D=E/A\) is \(N/A\) and
is nilpotent. The pure image in \(D^\dagger\) maps
onto \(N/A\) with central kernel, so is also nilpotent.
Thus its actual source map kills \(R_{\rm nil}(F)\)
and factors through \(\mathcal B_{\rm nil}\).
This verifies the required source factorization for
every modular auxiliary cover, with its torsion
relations already proved in \textup{(D.50.9)}.

\paragraph{Compactness with the complete conjugator coordinates.}
We prove \textup{(D.50.3)}; the modular proof uses the
same construction with the pure residual just described.
Work in the compact space \(\mathcal F^3\) of three
conjugators \((h_1,h_2,h_3)\). They specify
\(u_i=h_ix_ih_i^{-1}\).
Impose the relation \(u_1u_2u_3=1\) in every finite
quotient, generation by \(u_1,u_2\) in every finite
quotient, and the value of each \(h_i\) prescribed by
\(b_i\) in every finite nilpotent quotient.
Each is a closed condition, in fact a finite cylinder
condition. A finite collection of them is contained in
one actual joint finite quotient \(E_0\) of the source.
Put \(A_0=\gamma_\infty(E_0)\), \(D_0=E_0/A_0\), and
construct its actual central refinement \(D_0^\dagger\).
Every finite nilpotent constraint on \(E_0\) factors
through \(D_0\), and \(D_0^\dagger\) is itself
nilpotent. Define the new auxiliary epimorphism by
\[
 \mathcal F\xrightarrow{q_{\rm nil}}\mathcal F_{\rm nil}
       \xrightarrow{\alpha}\mathcal F_{\rm nil}
       \longrightarrow D_0^\dagger.
 \tag{D.50.19}
\]
Its three conjugators are the images of the specified
\(b_i\). The finite theorem gives an epimorphism to
\(E_0\) and actual conjugators there whose images in
\(D_0\), and hence in every prescribed nilpotent
factor, are those same values. Lift these three
elements along the original source map to \(E_0\).
They meet the entire finite collection of conditions.
This proves the finite intersection property.

The resulting compact intersection supplies three
actual \(h_i\), with \(q_{\rm nil}(h_i)=b_i\), satisfying
the full relation and every generation condition.
The universal property of the free profinite group gives
\(\Phi\) from \(x_1\mapsto u_1,x_2\mapsto u_2\).
Its image is compact. If it were proper, a finite
quotient would have proper image, contradicting the
generation condition. Thus \(\Phi\) is onto.
An onto endomorphism of a finitely generated profinite
group is injective: for each index there are only
finitely many open subgroups, preimage under an onto
map permutes them, and any nontrivial kernel element
can be excluded by one such open subgroup, contradicting
its periodicity under preimage. This proves that
\(\Phi\) is an automorphism.
Equation \textup{(D.50.2)} gives the exact commutative
map in \textup{(D.50.3)} on the two generators and
hence on the entire source.

For \(\mathcal B\), use the compact space \(\mathcal B^3\)
with \(x_1=s,x_2=r,x_3=\tau^{-1}\), and include the
\(\mathfrak S_3\) quotient in \(E_0\).
Every finite constraint on \(\mathcal B_{\rm nil}\)
has nilpotent pure image, so it kills
\(\gamma_\infty(q_{E_0}(F))\) and factors through the
chosen \(D_0\). Its central refinement has nilpotent
pure image as well. The same finite construction
therefore applies to the prescribed automorphism of
\(\mathcal B_{\rm nil}\).
At the compact limit, \(u_1,u_2\) are conjugates of
\(s,r\), so they satisfy \(u_1^2=u_2^3=1\) literally.
The universal property of the profinite free product
gives the endomorphism of \(\mathcal B\); finite
generation and the preceding argument make it an
automorphism. This proves the full modular statement.

If the specified conjugator of \(x_1\), or of \(\tau\),
is one in the retained quotient, its lifted conjugator
lies in that quotient's kernel. Conjugating the full
automorphism by its inverse fixes the corresponding
peripheral element literally and leaves the retained
automorphism unchanged. This is the asserted normalized
lift, with its exact inner correction.

\paragraph{The exact products of local sources.}
The canonical homomorphism
\[
            F/R_{\rm nil}(F)\longrightarrow\prod_p F^{(p)}
 \tag{D.50.20}
\]
is an isomorphism. Its kernel is exactly the displayed
nilpotent residual, since every finite nilpotent group
is the product of its Sylow groups. It is onto: the
actual source image in any finite product of finite
quotients of distinct \(F^{(p)}\) is subdirect, and
the successive Goursat common quotients are trivial
by coprime orders. Compactness then gives surjectivity
onto the entire product. This also proves the identical
description of \(\mathcal F_{\rm nil}\) in
\textup{(D.50.1)}.

Put \(B_p=\mathcal B/O^p(F)\), with its actual quotient
to \(\mathfrak S_3\). There is an exact isomorphism
\[
 \mathcal B_{\rm nil}
       \longrightarrow\mathop{\prod}_{p,\,\mathfrak S_3}B_p,
 \tag{D.50.21}
\]
where the right side consists of tuples with one common
\(\mathfrak S_3\)-value.
Injectivity is the kernel calculation in
\textup{(D.50.20)}. For surjectivity, choose one element
of \(\mathcal B_{\rm nil}\) lifting the common
\(\mathfrak S_3\)-value. Multiplying the tuple by its
inverse leaves an arbitrary tuple in \(\prod_pF^{(p)}\),
and \textup{(D.50.20)} supplies the required actual
kernel element. No splitting of the two-primary or
three-primary modular extension is used.

Let \(f_\ell\) be an actual pro-\(\ell\) gentle pair
as in the statement. D.44.26--32 construct its actual
automorphism of \(B_\ell\), with
\[
       s\longmapsto sf_\ell,\qquad
       r\longmapsto f_\ell^{-1}r,\qquad \tau\longmapsto\tau.
 \tag{D.50.22}
\]
It induces the identity on \(\mathfrak S_3\).
The positive torsion conjugators can be chosen in
\(F^{(\ell)}\): the old and new order-two subgroups,
and separately the old and new order-three subgroups,
are complements to the pro-\(\ell\) kernel in the
preimages of their same cyclic quotient groups.
Coprime complement conjugacy, at all finite quotients
and then by compactness of the conjugator sets, gives
those kernel conjugators. Here \(\ell\geq5\) is used.

Use \textup{(D.50.22)} in the \(\ell\)-component of
\textup{(D.50.21)} and the identity in every other
component. The common \(\mathfrak S_3\)-map is unchanged,
so this defines an actual continuous automorphism of
\(\mathcal B_{\rm nil}\).
The just chosen conjugators, together with identity
conjugators in the other components, form actual
elements of its pure kernel and prove peripherality.
It fixes \(\tau\), with specified conjugator one.
Its element \(f_0=s\alpha(s)\) has coordinates
\(f_\ell\) at \(\ell\) and one elsewhere, and lies in
\((F/R_{\rm nil}(F))'\): the closed derived subgroup
of a product is the product of its closed derived
subgroups, as is verified on every finite set of
coordinates and then by closure.

Apply the modular lifting theorem and its exact
\(\tau\)-normalization to obtain \(\Phi\), and set
\(f=s\Phi(s)\). Since \(\Phi\) retains the identity
map on the specified \(\mathfrak S_3\), this element
belongs to \(F\). The actual quotient
\(\mathcal B/F'\) has abelian pure subgroup and
\(\mathfrak S_3\) quotient; the map from \(F\) to
that pure subgroup factors through \(F/R_{\rm nil}(F)\).
As \(f_0\) has trivial pure abelianization, so does
\(f\). Thus \(f\in F'\), with literal image
\(f_\ell\).
The equations \(\Phi(\tau)=\tau\) and \(\Phi(s)^2=1\)
give the exact formulas
\[
 \Phi(s)=sf,\quad \Phi(r)=f^{-1}r,\quad
       \Phi(x)=x,\quad \Phi(y)=f^{-1}yf.
 \tag{D.50.23}
\]
The actual conjugation table D.44.8 sends the two
torsion relations \((sf)^2=(f^{-1}r)^3=1\) to the
two literal gentle equations
\[
           f(y,x)f(x,y)=1,\qquad
           f(z,x)f(y,z)f(x,y)=1,
                 \qquad z=(xy)^{-1}.
 \tag{D.50.24}
\]
These equalities take place in the same full free
profinite group with the words \(x,y\) in
\textup{(D.50.23)}.
For any \(b\in\mathcal B\), the restrictions to \(F\)
satisfy
\(\Phi\operatorname{Ad}_b\Phi^{-1}=
\operatorname{Ad}_{\Phi(b)b^{-1}}\operatorname{Ad}_b\).
The correcting element belongs to \(F\), because
\(\Phi\) induces the identity on \(\mathfrak S_3\).
This proves the claimed centralization of its actual
outer action and completes the local-to-full comparison.

\paragraph{The prosolvable statements and their exact comparison.}
The same proof with the stable derived residual
\(E^{(\infty)}\) gives the whole prosolvable assertion.
This residual is perfect, hence satisfies \([E,A]=A\);
its quotient is solvable, and a central extension of
a solvable group is solvable. Every auxiliary source
therefore factors through the full prosolvable quotient.
For \(\mathcal B\), use the derived residual of its
finite pure image. It is \(E\)-normal, and the entire
finite quotient is solvable because its pure quotient
and its \(\mathfrak S_3\) quotient are solvable.
The same compact conjugator proof applies verbatim.

For clarity the actual pure quotient is retained here
too. The prosolvable residual of \(F\) is characteristic
in \(F\) and hence normal in \(\mathcal B\).
The quotient of \(\mathcal B\) by it is prosolvable,
since every finite quotient has solvable pure subgroup
and solvable \(\mathfrak S_3\) quotient.
Conversely the image of \(F\) in the prosolvable
quotient of \(\mathcal B\) is prosolvable.
These two actual quotient maps prove that the two
residuals coincide and identify the pure kernel with
the full prosolvable quotient of \(F\).
All the normalizations above therefore retain the
specified prosolvable maps and literal old gentle pair.
\end{proof}

\paragraph{Proof and computation scope.}
The finite scalar covers, coefficient formula, every proper-image
subtraction, centrality, joint generation, and compactness maps in
this proof are algebraic arguments. The independent finite
certificates with prefix \texttt{bc\_scalar\_cover\_} check the
actual scalar phase, the quaternionic degree-two coefficient,
the noninvariant-block vanishings, and a nonfaithful determinant
case. Their role is to check these exact finite identities; the
all-source assertion follows from the full proof above.
The earlier perfect-residual route is recorded as the prosolvable
corollary. The lower central residual in
\textup{(D.50.18)} gives the stronger pronilpotent conclusion because
the constructed auxiliary kernels are central in their whole groups.
